\subsection{Calculation of invariant factors}

PROPOSITION 4. --- Let A be a principal ideal domain, let L be a free A-module with finite basis \((u_{j})\) \((1 \leqslant j \leqslant k)\) , let M be a submodule of L, let \((x_{i})\) be a system of generators of M and let \(A\alpha_{i}\) \((1 \leqslant i \leqslant n)\) be the invariant factors of M with respect to L. Then for \(1 \leqslant m \leqslant n\) the product \(\delta_{m} = \alpha_{1} \ldots \alpha_{m}\) is a gcd of the m-th order minors of the matrix whose columns are the coordinate vectors of the \(x_{i}\) with respect to the basis \((u_{j})\) .

By Th. 1 it is clear that \(M \subset \alpha_{1}L\) ; hence the coordinates of any element of M are multiples of \(\alpha_{1}\) . On the other hand, there exists an element x of M for which \(\alpha_{1}\) is a content in L. Expressing x as a linear combination of the \(x_{t}\) , it follows that \(\alpha_{1}\) is an element of the ideal generated by the coordinates of the \(x_{t}\) . As these are all multiples of \(\alpha_{1}\) , it follows that \(\alpha_{1}\) is indeed their gcd, and our assertion is proved in the case m = 1.

For general \(m\) , consider the \(m\) -th exterior power \(\Lambda\) M of M (III, p.~507). In the notation of Th. 1 there exists a basis \((a_i)\) for M where \(a_i = \alpha_i e_i\) ( \(1 \leqslant i \leqslant n\) ); thus \(\Lambda\) M has a basis consisting of the elements \(a_{i_1} \wedge \ldots \wedge a_{i_m}\) , as \((i_1, \ldots, i_m)\) runs through the set of strictly increasing sequences of \(m\) elements of \((1, n)\) . Now the elements \(e_{i_1} \wedge \ldots \wedge e_{i_m}\) belong to a basis \(B_m\) of \(\bigwedge^m L\) . Hence the canonical map from \(\bigwedge^m M\) into \(\bigwedge^m L\) is an isomorphism from \(\bigwedge^m M\) onto the submodule of \(\bigwedge^m L\) having as a basis the elements \((\alpha_{i_1} \ldots \alpha_{i_m}) e_{i_1} \wedge \ldots \wedge e_{i_m}\) , and we identify this submodule with \(\bigwedge^m M\) . Since \(\alpha_j\) is a multiple of \(\alpha_k\) for \(j \geqslant k\) , the elements \(\alpha_{i_1} \ldots \alpha_{i_m}\) are all multiples of \(\delta_m = \alpha_1 \ldots \alpha_m\) , and one of them is equal to \(\delta_m\) ; thus \(\delta_m\) is a gcd of the coordinates with respect to \(B_m\) of the elements of a system of generators of \(\bigwedge^m M\) . The first part of the argument shows that then \(\delta_m\) is a gcd of the set of coordinates of any system of generators of \(\bigwedge^m M\) , with respect to any basis of \(\bigwedge^m L\) . Taking the basis for \(\bigwedge^m L\) induced from the basis \((u_j)\) of \(L\) , and the system of generators for \(\bigwedge^m M\) consisting of exterior products of the \((x_i)\) , the expression for the coordinates of these products in terms of determinants (III, p.~528, Prop. 9) gives the stated result.

